% Generated draft: requires mathematical and source review.
\documentclass[11pt]{article}
\usepackage[margin=25mm]{geometry}
\usepackage{fontspec}
\setmainfont{Latin Modern Roman}
\usepackage{amsmath,amssymb,mathtools,unicode-math}
\setmathfont{Latin Modern Math}
\usepackage{enumitem,needspace,xurl,hyperref}
\setlength{\parindent}{0pt}
\setlength{\parskip}{5pt}
\setlength{\emergencystretch}{3em}
\newcommand{\N}{\mathbb{N}}
\newcommand{\Z}{\mathbb{Z}}
\newcommand{\Q}{\mathbb{Q}}
\newcommand{\R}{\mathbb{R}}
\newcommand{\C}{\mathbb{C}}
\newcommand{\F}{\mathbb{F}}
\newcommand{\A}{\mathbb{A}}
\newcommand{\PP}{\mathbb{P}}
\newcommand{\E}{\mathbb{E}}
\newcommand{\eps}{\varepsilon}
\newcommand{\dd}{\,\mathrm{d}}
\begin{document}
\section*{UnsolvedMath ID: 11000217}
\textbf{Source problem code:} catalogue entry
\subsection{Definitions and mathematical statement}
Let \(\mathcal H_1(\alpha)\) be the area-one stratum of abelian differentials with zero multiplicities \(\alpha\) summing to \(2g-2\), and \(\mathcal Q_1(\beta)\) the area-one quadratic-differential stratum with the prescribed orders \(\beta\), summing to \(4g-4\). If simple poles are allowed, their order is \(-1\); the source does not repeat that convention. The \(\operatorname{SL}(2,\R)\)-action postcomposes translation or half-translation charts by real linear maps of determinant one, preserving flat area.

The problem is to classify the ergodic measures for this action on every indicated stratum. In the standard probability-measure interpretation, determine all invariant Borel probability measures \(\mu\) such that every measurable invariant set has measure \(0\) or \(1\). The excerpt does not specify whether infinite locally finite measures or other measure classes are also intended; that changes the requested scope.

The source credits McMullen [McM3] with solving its Problems 9 and 10 for genus-two \emph{translation} surfaces, not for all genera or all quadratic strata. It also defines a minimal set to be a nonempty closed invariant subset with no proper nonempty closed invariant subset, and gives the \(\operatorname{SL}(2,\R)\)-orbit of a Veech surface as an example. Nonemptiness is the standard implicit qualification needed to exclude the empty subset.

The supplied record labels the problem solved/SOLVED-IN-LITERATURE, while its summary calls the Eskin--Mirzakhani--Mohammadi measure-classification progress partial and not the full picture. The summary names invariant ergodic-measure classification but gives no exact theorem or citation matching both strata and the chosen measure class. These status/scope differences remain recorded rather than independently resolved.
\subsection{Short English statement}
Classify invariant ergodic measures on area-one abelian and quadratic differential strata, with the measure class specified and the supplied genus-two and later classification reports kept distinct.
\subsection{Sources}
\begin{itemize}
\item \emph{Problems on Mapping Class Groups and Related Topics}, edited by Farb, Problem 10, ``Ergodic measures,'' supplied PDF page 260. \url{https://www.math.uchicago.edu/~farb/papers/mcgbook.pdf}. Recorded access 2026-07-29; extraction pdf-font-structure.
\item McMullen [McM3], source attribution for genus-two translation-surface Problems 9 and 10; full bibliography and Problem 9 are absent. The summary credits Eskin--Mirzakhani--Mohammadi with invariant ergodic-measure classification but supplies no complete reference or precise scope.
\item Local UnsolvedMath registry, \texttt{data/unsolvedmath/problems.json}, permanent ID \(11000217\), exact source code \texttt{catalogue entry}; solved/SOLVED-IN-LITERATURE fields coexist with a PARTIAL summary.
\end{itemize}
\end{document}
